\documentclass[../../main]{subfiles}

\begin{document}

\section{Funções entre Conjuntos Ordenados}
\label{sec:matematica:funcoes:conjuntos-ordenados}

\subsection{Funções Monótonas}

\begin{defn}[Função Monótona]
	\label{def:matematica:funcoes:ordem:monotona}

	Sejam

	\[
		(A,\le_A)
	\]

	e

	\[
		(B,\le_B)
	\]

	conjuntos parcialmente ordenados.

	Uma função

	\[
		f:A\to B
	\]

	é dita monótona (ou crescente) quando

	\[
		x\le_A y
		\implies
		f(x)\le_B f(y)
	\]

	para quaisquer

	\[
		x,y\in A.
	\]

\end{defn}

\begin{obs}

	Também se utiliza a expressão
	\emph{preservadora de ordem}.

\end{obs}

\begin{ex}

	A função

	\[
		f:\mathbb R\to\mathbb R,
		\qquad
		f(x)=2x+1
	\]

	é monótona.

\end{ex}

\begin{prop}
	\label{prop:matematica:funcoes:ordem:identidade}

	A função identidade de qualquer conjunto ordenado é monótona.

\end{prop}

\begin{proof}

	Se

	\[
		x\le y,
	\]

	então

	\[
		\operatorname{id}(x)
		=
		x
		\le
		y
		=
		\operatorname{id}(y).
	\]

\end{proof}

\begin{prop}
	\label{prop:matematica:funcoes:ordem:composicao}

	A composição de funções monótonas é monótona.

\end{prop}

\begin{proof}

	Se

	\[
		x\le_A y,
	\]

	então

	\[
		f(x)\le_B f(y).
	\]

	Como \(g\) é monótona,

	\[
		g(f(x))
		\le_C
		g(f(y)).
	\]

	Portanto

	\[
		(g\circ f)(x)
		\le_C
		(g\circ f)(y).
	\]

\end{proof}

\subsection{Funções Antitônicas}

\begin{defn}[Função Antitônica]
	\label{def:matematica:funcoes:ordem:antitona}

	Sejam

	\[
		(A,\le_A)
	\]

	e

	\[
		(B,\le_B).
	\]

	Uma função

	\[
		f:A\to B
	\]

	é dita antitônica (ou decrescente) quando

	\[
		x\le_A y
		\implies
		f(y)\le_B f(x).
	\]

\end{defn}

\begin{ex}

	A função

	\[
		f:\mathbb R\to\mathbb R,
		\qquad
		f(x)=-x
	\]

	é antitônica.

\end{ex}

\begin{prop}
	\label{prop:matematica:funcoes:ordem:antitona-composicao}

	A composição de duas funções antitônicas é monótona.

\end{prop}

\begin{proof}

	A inversão da ordem ocorre duas vezes.

\end{proof}

\begin{prop}
	\label{prop:matematica:funcoes:ordem:monotona-antitona}

	A composição de uma função monótona com uma função antitônica é
	antitônica.

\end{prop}

\begin{proof}

	Segue diretamente das definições.

\end{proof}

\subsection{Caracterização por Relações de Ordem}

\begin{prop}
	\label{prop:matematica:funcoes:ordem:caracterizacao}

	Uma função

	\[
		f:A\to B
	\]

	é monótona se, e somente se,

	\[
		(x,y)\in\le_A
		\implies
		(f(x),f(y))\in\le_B.
	\]

\end{prop}

\begin{proof}

	Trata-se apenas da reescrita da definição usando relações.

\end{proof}

\subsection{Imersões de Ordem}

\begin{defn}[Imersão de Ordem]
	\label{def:matematica:funcoes:ordem:imersao}

	Seja

	\[
		f:A\to B.
	\]

	Diz-se que \(f\) é uma imersão de ordem quando

	\[
		x\le_A y
		\iff
		f(x)\le_B f(y).
	\]

\end{defn}

\begin{prop}
	\label{prop:matematica:funcoes:ordem:imersao-injetiva}

	Toda imersão de ordem é injetiva.

\end{prop}

\begin{proof}

	Suponha

	\[
		f(x)=f(y).
	\]

	Então

	\[
		f(x)\le f(y)
	\]

	e

	\[
		f(y)\le f(x).
	\]

	Pela definição de imersão,

	\[
		x\le y
	\]

	e

	\[
		y\le x.
	\]

	Pela antissimetria,

	\[
		x=y.
	\]

\end{proof}

\subsection{Isomorfismos de Ordem}

\begin{defn}[Isomorfismo de Ordem]
	\label{def:matematica:funcoes:ordem:isomorfismo}

	Um isomorfismo de ordem é uma bijeção

	\[
		f:A\to B
	\]

	tal que

	\[
		x\le_A y
		\iff
		f(x)\le_B f(y).
	\]

\end{defn}

\begin{prop}
	\label{prop:matematica:funcoes:ordem:isomorfismo-inversa}

	A inversa de um isomorfismo de ordem é um isomorfismo de ordem.

\end{prop}

\begin{proof}

	Seja

	\[
		y_1\le_B y_2.
	\]

	Escreva

	\[
		y_1=f(x_1),
		\qquad
		y_2=f(x_2).
	\]

	Como \(f\) preserva e reflete ordem,

	\[
		x_1\le_A x_2.
	\]

	Portanto

	\[
		f^{-1}(y_1)
		\le_A
		f^{-1}(y_2).
	\]

\end{proof}

\subsection{Preservação de Supremos}

\begin{defn}[Preservação de Supremo]
	\label{def:matematica:funcoes:ordem:supremo}

	Seja

	\[
		f:A\to B.
	\]

	Diz-se que \(f\) preserva supremos quando

	\[
		f(\sup X)
		=
		\sup f[X]
	\]

	para todo subconjunto

	\[
		X\subseteq A
	\]

	cujo supremo exista.

\end{defn}

\begin{obs}

	Nem toda função monótona preserva supremos.

\end{obs}

\begin{ex}

	Em reticulados completos, muitas funções importantes preservam
	supremos arbitrários.

\end{ex}

\subsection{Espaço das Funções Monótonas}

\begin{defn}[Espaço das Funções Monótonas]
	\label{def:matematica:funcoes:ordem:espaco-monotono}

	Sejam

	\[
		(A,\le_A)
	\]

	e

	\[
		(B,\le_B).
	\]

	Denota-se por

	\[
		\operatorname{Mon}(A,B)
	\]

	o conjunto de todas as funções monótonas de \(A\) em \(B\).

\end{defn}

\begin{prop}
	\label{prop:matematica:funcoes:ordem:ordem-ponto-a-ponto}

	A relação

	\[
		f\le g
		\iff
		(\forall x\in A)
		\;
		f(x)\le_B g(x)
	\]

	define uma ordem parcial em

	\[
		\operatorname{Mon}(A,B).
	\]

\end{prop}

\begin{proof}

	Reflexividade, antissimetria e transitividade seguem da ordem
	ponto a ponto em \(B\).

\end{proof}

\end{document}
